\documentclass[a4paper,12pt]{article}
\usepackage{graphicx}
\usepackage[left=2.5cm, right=2.5cm, top=3cm, bottom=3cm]{geometry}
\usepackage{amsmath, amsthm, amssymb}
\usepackage[left=2.5cm, right=2.5cm, top=3cm, bottom=3cm]{geometry}
\usepackage[hidelinks]{hyperref}
\usepackage[spanish]{babel}
\usepackage{cite}
\usepackage[usenames]{color}
\usepackage{url}

\bibliographystyle{plain}
\begin{document}
\color{white}
\pagecolor{black}
    \ttfamily
    \title{\huge \textbf{H.U.L.K}}
    \author{Diego Manuel Viera Martínez}
    \date{octubre , 2023}
    \maketitle
    \begin{center}
        \includegraphics[scale=0.7]{Pictures/matcom.jpg}
        \label{fig:logo}
    \end{center}
    \begin{center}
        \large Proyecto II de programación
    \end{center}

    \newpage 
    
    \section{Expresiones H.U.L.K}\label{sec:Expresiones}
    \subsection{Print:}\label{sub:printExpression}
    \begin{center}
        Una de las instruccioines más simples en H.U.L.K (print),
    escirbe en consola el valorde la expresión que pongas entre
    paréntesis después del "print", en caso de no poner ninguna expresión
    pone una linea en blanco.
    \end{center}
    \textcolor{green}{$>$ print("Hello World !");}\\
    \textcolor{green}{Hello World !}
    

    \subsection{Operaciones Aritméticas:}\label{sub:aritmeticsOperations}
    \begin{flushleft}
        H.U.L.K cuenta con los operadores aritméticos: + (suma),  - (resta), 
        * (multiplicación), / (división), \% (módulo) y potencia.
    \end{flushleft}
    \textcolor{green}{$>$ print(((1 + 2) * 4) / 2);}
    
    \textcolor{green}{6}

    \subsection{Funciones matemáticas básicas:}\label{sub:mathExpressions}
    \begin{flushleft}
        Contiene funciones matemáticas básicas para operaciones:
        
        \begin{enumerate}
            \item `sqrt(<value>)` calcula la raíz cuadrada de un valor.
            \item `sin(<angle>)` calcual en seno de un ángulo en radianes.
            \item `cos(<angle>)` calcual en coseno de un ángulo en radianes.
            \item exp(<value>)` calcula el valor de `e` elevado a un valor.
            \item `log(<base>, <value>)` calcula el logaritmo de un valor en una base dada.
            \item `rand()` devuelve un número uniforme aleatorio entre 0 y 1 (ambos inclusive).
        \end{enumerate}
    \end{flushleft}

    \textcolor{green}{$>$ print(sin(2 * PI) * 2 + cos(3 * PI / log(4, 64)));}
    
    \subsection{Funciones:}\label{sub:functions}
    \begin{flushleft}
        Para declarar funciones hay se debe seguir una sintaxis:

        function $<$nombre de la función$>$(parámetros) =>  $<$cuerpo de la función$>$

    \end{flushleft}
    Ejemplo:

    \textcolor{green}{$>$ function tan(x) => sin(x) / cos(x);}

    \newpage

    \subsection{Expresiones let-in:}\label{sub:letin}
    \begin{flushleft}
        En H.U.L.K es posible declarar variables usando la expresión let-in, que funciona de la siguiente forma:
    \end{flushleft}

    \textcolor{green}{$>$ let x = PI/2 in print(tan(x));}

    \begin{flushleft}
        En general, una expresión let-in consta de una o más declaraciones de variables, y un cuerpo, que puede ser cualquier expresión donde además se pueden utilizar las variables declaradas en el let. Fuera de una expresión let-in las variables dejan de existir.
    \end{flushleft}

    \begin{flushleft}
        El valor de retorno de una expresión let-in es el valor de retorno del cuerpo, por lo que es posible hacer:
    \end{flushleft}
    
    \textcolor{green}{$>$ print(7 + (let x = 2 in x * x));}
    
    \subsection{Condicionales}\label{sub:conditionals}
        \begin{flushleft}
            Las condiciones en HULK se implementan con la expresión if-else, que recibe una expresión booleana entre paréntesis, y dos expresiones para el cuerpo del if y el else respectivamente. Siempre deben incluirse ambas partes:
        \end{flushleft}
        
        \textcolor{green}{$>$ let a = 42 in if (a \% 2 == 0) print($"$Even$"$) else print($"$odd$"$);}
        
        \begin{flushleft}
            Como if-else es una expresión, se puede usar dentro de otra expresión
        \end{flushleft}

        \textcolor{green}{$>$ let a = 42 in print(if (a \% 2 == 0) $"$even$"$ else $"$odd$"$);}

    \subsubsection{Recursión}\label{sub:recursive}
    \begin{flushleft}
        H.U.L.K soporta declaraciones de funciones recursivas:
    \end{flushleft}

    \textcolor{green}{$>$ function fib(n) => if (n > 1) fib(n-1) + fib(n-2) else 1;}

    \subsection{Errores}\label{sub:errors}
    \begin{flushleft}
        Presenta 3 tipos de errores:
        \begin{enumerate}
            \item Lexical Errors: Errores que se producen por la presencia de tokens inválidos.\\
            \textcolor{green}{$>$ let 14a = 5 in print(14a);}\\
            \textcolor{green}{! LEXICAL ERROR: `14a` is not valid token.}
            \\
            \item Syntax Errors: Errores que se producen por expresiones mal formadas.\\
            \textcolor{green}{$>$ let a = 5 in print(a;}\\
            \textcolor{green}{! SYNTAX ERROR: Missing closing parenthesis after `a`.}\\

            \textcolor{green}{$>$ let a = 5 inn print(a);}\\
            \textcolor{green}{! SYNTAX ERROR: Invalid token `inn` in `let-in` expression.}\\

            \textcolor{green}{$>$ let a = in print(a);}\\
            \textcolor{green}{! SYNTAX ERROR: Missing expression in `let-in` after variable `a`.}\\

            \newpage
            \item Semantic Errors: Errores que se producen por el uso incorrecto de los tipos y argumentos.\\
            \textcolor{green}{$>$ let a = "hello world" in print(a + 5);}\\
            \textcolor{green}{! SEMANTIC ERROR: Operator `+` cannot be used between `string` and `number`.}\\

            \textcolor{green}{$>$ print(fib("hello world"));}\\
            \textcolor{green}{! SEMANTIC ERROR: Function `fib` receives `number`, not `string`.}\\
            
            \textcolor{green}{$>$  print(fib(4,3));}\\
            \textcolor{green}{! SEMANTIC ERROR: Function `fib` receives 1 argument(s), but 2 were given.}\\
        \end{enumerate}
    \end{flushleft}

    \subsection{¿ Cómo iniciar y detener ?}\label{sub:final}
    \begin{flushleft}
        Para inicializar H.U.L.K debes abrir su carpeta contenedora en consola y usar el comando(dotnet run)
        y listo puedes comenzar a usarlo. Para cerrar el porgrama escribe en el input (stop hulk) y regresará a la consola de su PC.
    \end{flushleft}
\end{document}
